% ==============================================================================
% How to Design Programs (Second Edition) - 紧凑双栏中文排版宏包与样式配置
% ==============================================================================

% 字号配置：全局 9pt 紧凑版面
\usepackage[fontsize=9pt]{scrextend}

% 页面几何与边距（左右边距减半至 0.8cm，可用版心大幅扩展）
\usepackage[a4paper, top=1.8cm, bottom=1.8cm, left=0.8cm, right=0.8cm, headheight=13pt, footskip=18pt, columnsep=6.5mm]{geometry}
\setlength{\columnseprule}{0pt} % 去除双栏中线

% 数学与符号
\usepackage{amsmath, amssymb}
\usepackage{newunicodechar}
\newunicodechar{λ}{\ensuremath{\lambda}}
\newunicodechar{ε}{\ensuremath{\varepsilon}}
\newunicodechar{⊕}{\ensuremath{\oplus}}

% 颜色定义
\usepackage[x11names, svgnames, dvipsnames]{xcolor}
\definecolor{NavyTheme}{RGB}{24, 43, 73}        % 主题深蓝
\definecolor{AccentGreen}{RGB}{26, 127, 86}     % 配方绿
\definecolor{AccentPurple}{RGB}{103, 58, 183}   % 练习紫
\definecolor{AccentAmber}{RGB}{217, 119, 6}     % 注意橙
\definecolor{CodeBg}{RGB}{247, 248, 250}        % 代码背景浅灰
\definecolor{CodeFrame}{RGB}{220, 224, 230}     % 代码边框
\definecolor{ReplBg}{RGB}{243, 244, 246}        % REPL背景
\definecolor{TextDark}{RGB}{30, 41, 59}         % 正文字体深色
\definecolor{StepHl}{RGB}{216, 221, 230}        % 单步求值高亮底色
\definecolor{StepHlDeep}{RGB}{190, 197, 210}    % 嵌套作用域深高亮

% 图形与自适应双栏限宽
\usepackage{graphicx}
\usepackage[export]{adjustbox}
\graphicspath{{images/}}

% 拦截所有图片，保证在单栏内绝对不超出栏宽，保持纵横比
\let\oldincludegraphics\includegraphics
\RenewDocumentCommand{\includegraphics}{O{} m}{%
  \oldincludegraphics[#1, max width=\linewidth, keepaspectratio]{#2}%
}

\usepackage{caption}
\captionsetup{font=footnotesize, labelfont={bf, color=NavyTheme}, labelsep=quad, skip=3pt}

% 段落与列表紧凑化
\linespread{1.12}
\setlength{\parskip}{1.5pt plus 1pt minus 0.5pt}
\setlength{\parindent}{2\ccwd}
\usepackage{enumitem}
\setlist{itemsep=1.5pt, parsep=0.5pt, topsep=2.5pt, partopsep=0pt, leftmargin=1.5em}

% 页眉页脚
\usepackage{fancyhdr}
\pagestyle{fancy}
\fancyhf{}
\fancyhead[LE]{\small\color{gray}\nouppercase{\leftmark}}
\fancyhead[RO]{\small\color{gray}\nouppercase{\rightmark}}
\fancyfoot[C]{\small\thepage}
\renewcommand{\headrulewidth}{0.3pt}
\renewcommand{\footrulewidth}{0pt}

\fancypagestyle{plain}{
  \fancyhf{}
  \fancyfoot[C]{\small\thepage}
  \renewcommand{\headrulewidth}{0pt}
}

% 紧凑标题定制
\usepackage{titlesec}
\titleformat{\chapter}[block]
  {\normalfont\LARGE\bfseries\color{NavyTheme}}
  {\color{NavyTheme!70}\thechapter\quad}
  {0pt}
  {\LARGE}
\titlespacing*{\chapter}{0pt}{-10pt}{10pt}

\titleformat{\section}
  {\normalfont\large\bfseries\color{NavyTheme}}
  {\thesection}{0.6em}{}
\titlespacing*{\section}{0pt}{9pt}{3pt}

\titleformat{\subsection}
  {\normalfont\normalsize\bfseries\color{NavyTheme!85}}
  {\thesubsection}{0.5em}{}
\titlespacing*{\subsection}{0pt}{6pt}{2pt}

\titleformat{\subsubsection}
  {\normalfont\small\bfseries\color{NavyTheme!80}}
  {\thesubsubsection}{0.5em}{}
\titlespacing*{\subsubsection}{0pt}{4pt}{2pt}

% 代码高亮配置 (Listings & Stepper)
\usepackage{listings}
\usepackage{fvextra}
\usepackage{booktabs}
\usepackage{array}
\lstdefinelanguage{BSL}{
  keywords={define, define-struct, check-expect, check-within, check-error, check-satisfied,
            cond, else, if, and, or, not, require, local, let, lambda, empty, match},
  keywordstyle=\bfseries\color{NavyTheme},
  basicstyle=\footnotesize\ttfamily,
  stringstyle=\color{AccentGreen},
  commentstyle=\color{gray}\itshape,
  showstringspaces=false,
  breaklines=true,
  tabsize=2,
  columns=fullflexible,
  numbers=none,
  morecomment=[l]{;},
  morestring=[b]",
  literate={{\#true}}{{\textcolor{AccentAmber}{\#true}}}5
           {{\#false}}{{\textcolor{AccentAmber}{\#false}}}6
           {'()}{{\textcolor{NavyTheme}{\textquotesingle()}}}3
           {λ}{{$\lambda$}}1
           {ε}{{$\varepsilon$}}1
}

\lstset{
  language=BSL,
  basicstyle=\footnotesize\ttfamily,
  backgroundcolor=\color{CodeBg},
  frame=single,
  rulecolor=\color{CodeFrame},
  frameround=tttt,
  framesep=3pt,
  xleftmargin=3pt,
  xrightmargin=3pt,
  keepspaces=true,
  breaklines=true,
  literate={{\#true}}{{\textcolor{AccentAmber}{\#true}}}5
           {{\#false}}{{\textcolor{AccentAmber}{\#false}}}6
           {'()}{{\textcolor{NavyTheme}{\textquotesingle()}}}3
           {λ}{{$\lambda$}}1
           {ε}{{$\varepsilon$}}1
}

% 行内代码样式
% 通过 xeCJK 边界/空格节点让 \colorbox 与前后中文及空格保持对称间距
\ExplSyntaxOn
\NewDocumentCommand{\code}{m}
  {
    \ifmmode
      \textnormal{\setlength{\fboxsep}{2pt}\colorbox{CodeBg}{\footnotesize\ttfamily\textcolor{NavyTheme}{#1}}}
    \else
      \leavevmode
      \xeCJK_Boundary_and_Default:
      \group_begin:
        \setlength{\fboxsep}{2pt}
        \colorbox{CodeBg}{\footnotesize\ttfamily\textcolor{NavyTheme}{#1}}
      \group_end:
      \peek_catcode:NTF \c_space_token
        { \xeCJK_make_space_node: }
        { { \xeCJK_make_node:n { default } } }
    \fi
  }
\ExplSyntaxOff

% 彩色方框定义 (tcolorbox) - 紧凑双栏优化
\usepackage[most]{tcolorbox}

% 1. 设计配方环境 (Design Recipe)
\newtcolorbox{recipebox}[1][]{
  enhanced,
  breakable,
  colback=AccentGreen!4!white,
  colframe=AccentGreen!85!black,
  fonttitle=\small\bfseries\sffamily,
  title={【设计配方】 #1},
  arc=1.5mm,
  boxrule=0.7pt,
  boxsep=2pt,
  left=6pt, right=6pt, top=4pt, bottom=4pt,
  before skip=5pt, after skip=5pt
}

% 2. 练习环境 (Exercise)
\newtcbtheorem{exercise}{练习}{
  enhanced,
  breakable,
  colback=AccentPurple!3!white,
  colframe=AccentPurple!80!black,
  fonttitle=\small\bfseries\sffamily,
  arc=1.5mm,
  boxrule=0.7pt,
  boxsep=2pt,
  left=6pt, right=6pt, top=4pt, bottom=4pt,
  before skip=5pt, after skip=5pt
}{ex}

% 3. 提示与核心要点环境 (Note / Takeaway)
\newtcolorbox{notebox}[1][核心要点]{
  enhanced,
  breakable,
  colback=AccentAmber!6!white,
  colframe=AccentAmber!90!black,
  fonttitle=\small\bfseries\sffamily,
  title={【要点】 #1},
  arc=1.5mm,
  boxrule=0.7pt,
  boxsep=2pt,
  left=6pt, right=6pt, top=4pt, bottom=4pt,
  before skip=5pt, after skip=5pt
}

% 4. 交互式 REPL 会话环境 (DrRacket Interactions)
\newtcblisting{replbox}{
  enhanced,
  breakable,
  listing only,
  listing options={
    language=BSL,
    basicstyle=\footnotesize\ttfamily,
    breaklines=true,
    columns=fullflexible,
    keepspaces=true
  },
  colback=ReplBg,
  colframe=TextDark!80!white,
  coltitle=white,
  fonttitle=\footnotesize\ttfamily\bfseries,
  title={DrRacket 交互窗口},
  arc=1.2mm,
  boxrule=0.7pt,
  boxsep=2pt,
  left=6pt, right=6pt, top=3pt, bottom=3pt,
  before skip=5pt, after skip=5pt
}

% 5. 术语与旁注标注
\newcommand{\term}[1]{\textbf{\textcolor{NavyTheme}{#1}}}
\newcommand{\enquote}[1]{“#1”}

% 6. 单步求值/语法/对比代码块环境与高亮宏
\newtcolorbox{stepperbox}[1][]{
  enhanced,
  breakable,
  colback=CodeBg,
  colframe=CodeFrame,
  arc=1.2mm,
  boxrule=0.5pt,
  boxsep=2pt,
  left=4pt, right=4pt, top=3pt, bottom=3pt,
  before skip=4pt, after skip=4pt,
  #1
}
\newcommand{\hl}[1]{{\setlength{\fboxsep}{0.6pt}\colorbox{StepHl}{\vphantom{Ag}#1}}}
\newcommand{\hld}[1]{{\setlength{\fboxsep}{0.6pt}\colorbox{StepHlDeep}{\vphantom{Ag}#1}}}
\newcommand{\inlinesprite}[2][1.5em]{\raisebox{-0.2\height}{\oldincludegraphics[height=#1, keepaspectratio]{#2}}}

% 超链接配置
\usepackage{hyperref}
\usepackage{bookmark}
\hypersetup{
  colorlinks=true,
  linkcolor=NavyTheme,
  citecolor=AccentGreen,
  urlcolor=NavyTheme!80!blue,
  pdfborder={0 0 0},
  bookmarksnumbered=true,
  pdftitle={程序设计方法 (How to Design Programs, Second Edition)},
  pdfauthor={Matthias Felleisen, Robert Bruce Findler, Matthew Flatt, Shriram Krishnamurthi / Minsecrus 译}
}


% 章节与插曲图表计数器与编号自动规范化 (图 <章号>.<图号>)
\let\origchapter\chapter
\RenewDocumentCommand{\chapter}{s o m}{%
  \IfBooleanTF{#1}{%
    \origchapter*{#3}%
  }{%
    \IfNoValueTF{#2}{%
      \origchapter{#3}%
    }{%
      \origchapter[#2]{#3}%
    }%
    \ifnum\value{chapter}>0
      \renewcommand{\thefigure}{\arabic{chapter}.\arabic{figure}}%
      \setcounter{figure}{0}%
      \renewcommand{\thetable}{\arabic{chapter}.\arabic{table}}%
      \setcounter{table}{0}%
    \fi
  }%
}

% 插曲独立单元定义（不占用正文章节序号，全书 34 章与原书严格 1:1 对齐）
\newcounter{intermezzocnt}
\newcommand{\intermezzo}[1]{%
  \stepcounter{intermezzocnt}%
  \clearpage
  \origchapter*{#1}%
  \addcontentsline{toc}{chapter}{#1}%
  \markboth{#1}{#1}%
  \renewcommand{\thefigure}{插曲\arabic{intermezzocnt}.\arabic{figure}}%
  \setcounter{figure}{0}%
  \renewcommand{\thetable}{插曲\arabic{intermezzocnt}.\arabic{table}}%
  \setcounter{table}{0}%
}
\newcommand{\intermezzosection}[1]{%
  \section*{#1}%
  \addcontentsline{toc}{section}{#1}%
}
